%%%%%%%%%%%%%%%%%%%%%%%%%%%%%%%%%%%%%%%%%
% KaoNotes multilingual and code example
%%%%%%%%%%%%%%%%%%%%%%%%%%%%%%%%%%%%%%%%%

\documentclass[
	a4paper,
	fontsize=10pt,
	open=any,
	titlepage=false,
	numbers=noenddot,
]{kaobook}

\ifxetexorluatex
	\usepackage{polyglossia}
	\setmainlanguage{english}
	\usepackage{xeCJK}
	\setCJKmainfont{Noto Serif TC}
	\setCJKsansfont{Noto Sans TC}
	\setCJKmonofont{Noto Sans TC}
	\newCJKfontfamily\zhHans{LXGW Neo ZhiSong}[AutoFakeBold=2]
\else
	\usepackage[english]{babel}
\fi

\usepackage[framed=true]{kaotheorems}
\usepackage{kaorefs}

\title{KaoNotes Usage Examples}
\subtitle{English, Traditional Chinese, Simplified Chinese, and Code}
\author{J Song (Alex)}
\date{}

\begin{document}

\maketitle

\chapter{Multilingual and Code Examples}

\section{English mathematics}

For a matrix \(A\in\mathbb{R}^{m\times n}\) and a vector
\(\mathbf{x}\in\mathbb{R}^{n}\), the product \(A\mathbf{x}\) is a linear
combination of the columns of \(A\).

\begin{definition}
The column space of \(A\) is the set of all vectors that can be written as
\(A\mathbf{x}\).
\end{definition}

\section{繁體中文範例}

若矩陣 \(A\) 的行向量線性獨立，則方程 \(A\mathbf{x}=\mathbf{b}\)
對每個可達的目標向量都有唯一表示。KaoNotes 可用於整理數學推導、
演算法筆記、課堂內容與研究紀錄。

\begin{remark}
中英文可以在同一段落中混合排版，例如：矩陣乘法
\emph{matrix multiplication} 與向量空間 \emph{vector space}。
\end{remark}

\begingroup
\zhHans

\section{\zhHans 简体中文示例}

矩阵可以表示线性变换。通过列空间的观点，求解
\(A\mathbf{x}=\mathbf{b}\) 等价于判断目标向量 \(\mathbf{b}\)
能否由矩阵 \(A\) 的列向量线性组合得到。

\begin{kaobox}[frametitle={\zhHans 书写范围}]
可用于数学公式、证明、计算机科学笔记、算法说明、代码示例和研究记录。
\end{kaobox}

\endgroup

\section{Python code example}

Use the standard \texttt{lstlisting} environment for source code.

\begin{lstlisting}[language=Python, caption={Matrix-vector multiplication}]
def matvec(matrix, vector):
    """Return the matrix-vector product."""
    return [
        sum(value * vector[j] for j, value in enumerate(row))
        for row in matrix
    ]

A = [[2, 5], [1, 3]]
x = [1, 2]
print(matvec(A, x))  # [12, 7]
\end{lstlisting}

\end{document}
